\documentclass[11pt]{article}
\usepackage{mathblatt}
% Präambel des Setzers (werkzeuge/setzer-praeambel.tex, wortgleich); dazu \gteil (Muster T2U)
\makeatletter
\setlength{\mbpfbe}{0mm}% keine Punkte (pruefung.md § 5: nur beim Prüfstein)
\newlength{\gEin}\setlength{\gEin}{0pt}
\let\gSchrift\relax
\renewcommand{\pfkreuzzeile}[1]{\par\noindent\hangindent3.2em\hangafter1\hspace*{1.6em}$\square$\ #1\par}
\newcommand{\gkreuz}[1]{$\square$\ #1\hspace{2em}}
\newcommand{\gkreuzzeile}[2]{\par\noindent\hangindent2.2em\hangafter1\mbox{#1\,$\square$}\ #2\par}% Buchstabe, Option: eine je Zeile (neben Skizze)
% links, was man ansieht, rechts, was man tut; Spaltengrenze überall an derselben Stelle (gemeinsam.md „Satz“)
\newsavebox{\gbox}\newlength{\gLinks}
% v0.10 (Lehrer 13.10.): keine starre Halbierung – die linke Spalte ist so breit wie die Figur (+ 6 mm), die
% Antwort steht direkt daneben; ist die Figur breiter als 0,62 Zeile, steht die Antwort darunter
\newcommand{\gzwei}[2]{\par\smallskip\noindent\sbox{\gbox}{#1}%
  \setlength{\gLinks}{\dimexpr\wd\gbox+6mm\relax}%
  \ifdim\gLinks>0.62\linewidth
    \usebox{\gbox}\par\smallskip\noindent\begin{minipage}{\linewidth}#2\end{minipage}%
  \else
    \begin{minipage}[t]{\gLinks}\vspace{0pt}\usebox{\gbox}\end{minipage}%
    \begin{minipage}[t]{\dimexpr\linewidth-\gLinks\relax}\vspace{0pt}\raggedright #2\end{minipage}%
  \fi\par}
% v0.9 (12.10., pruefung.md § 5 „Marke mit B, Bedingungen rechts“): drittes Argument von pfaufg = Bedingungen
% der Prüfung (v0.10: \pfohnetr = TR durchgestrichen, „FOR“), klein grau im rechten Rand auf der Grundlinie der Nummer, Zeilen mit \\
\newcommand{\gbedrand}[2]{\ifblank{#2}{}{\rlap{\hspace*{\dimexpr#1+2mm\relax}%
  \parbox[t]{15mm}{\raggedright\scriptsize\leavevmode\color{mbgrau}\strut#2}}}}
% Teil eines Bündels (a, b …): Marke links im Rand, Bedingungen rechts im Rand, auf der Zeile des Teils
\newcommand{\gteilrand}[2]{\llap{\mbpfrandmarke{#1}\hspace*{\mbpfnr}}\gbedrand{\linewidth}{#2}}
% Teil eines Bündels: Buchstabe in eigener Spalte, Inhalt eingerückt (\gEin = Einzug, damit die Spaltengrenze von
% \gzwei an derselben Stelle bleibt)
\newenvironment{gbteil}[3]{\par\smallskip\noindent\gteilrand{#2}{#3}\makebox[\mbpfteil][l]{\textbf{#1}}%
  \begin{minipage}[t]{\dimexpr\linewidth-\mbpfteil\relax}\setlength{\gEin}{\mbpfteil}\raggedright\strut\ignorespaces}%
 {\end{minipage}\par}
% Grafik höchstens so breit wie die Spalte (Bündel nebeneinander, Koordinatensystem ohne Karo)
\newsavebox{\gpbox}
\newcommand{\gpassend}[1]{\sbox{\gpbox}{#1}\ifdim\wd\gpbox>\linewidth\typeout{GPASSEND: verkleinert \the\wd\gpbox\space auf \the\linewidth}\resizebox{\linewidth}{!}{\usebox{\gpbox}}\else\usebox{\gpbox}\fi}
% Skizzenfeld (v0.10, Lehrer 13.10.): Karo ab dem linken Rand über die Zeile, „Skizze“ klein grau oben links im Feld
% (wie \pfskizzenfeld der Vorlage); der rechte Rand bleibt den Bedingungen
\newcommand{\gskizze}[1]{\par\smallskip\noindent\begin{tikzpicture}
  \clip (0,0) rectangle ({\linewidth-1mm},{#1*10mm});
  \draw[black!16,line width=0.3pt,step=5mm] (0,0) grid ({\linewidth},{#1*10mm});
  \node[anchor=north west,inner sep=2pt,font=\scriptsize,text=mbgrau,fill=white] at (0,{#1*10mm}) {Skizze};
  \end{tikzpicture}\par}
\RenewDocumentEnvironment{pfaufg}{m m m}{%
  \begin{lrbox}{\mbaufgabebox}\begin{minipage}[t]{\linewidth}%
  \gSchrift\noindent\hspace*{\gEin}\mbpfrandmarke{#1}\gbedrand{\linewidth-\gEin}{#3}%
  \begin{minipage}[t]{\mbpfnr}\strut\textbf{#2}\end{minipage}%
  \begin{minipage}[t]{\dimexpr\linewidth-\gEin-\mbpfnr\relax}\raggedright\strut\ignorespaces}%
 {\end{minipage}%
  \end{minipage}\end{lrbox}%
  \setlength{\mbaufgabehoehe}{\dimexpr\ht\mbaufgabebox+\dp\mbaufgabebox\relax}%
  \ifdefempty{\gkopfnach}{\par\addvspace{5pt}\Needspace*{\mbaufgabehoehe}}{%
    % Stufenkopf, der mit dieser (ersten) Aufgabe zusammenbleibt: Platz für beide verlangen, dann den Kopf setzen
    % (12.10., Quadratische: ein Kopf stand allein unten, die Aufgabe mit Koordinatensystem auf der nächsten Seite)
    \par\Needspace*{\dimexpr\mbaufgabehoehe+40pt\relax}\gkopfnach\gdef\gkopfnach{}\par\addvspace{5pt}}%
  \noindent\usebox{\mbaufgabebox}\par\addvspace{5pt}}
% Stufenkopf vormerken: \gkopfvor{<Kopf>} – die nächste pfaufg setzt ihn, sobald beide zusammen Platz haben
\def\gkopfnach{}
\newcommand{\gkopfvor}[1]{\gdef\gkopfnach{#1}}
% Teilaufgabe mit Skizze links in gzwei: Buchstabe in eigener Spalte, darunter Text und Skizze
\newcommand{\gteil}[2]{\makebox[\mbpfteil][l]{\textbf{#1}}\parbox[t]{\dimexpr0.5\textwidth-\gEin-\mbpfnr-5mm-\mbpfteil\relax}{\raggedright #2}}
% Stufenkopf der Vorlage (mathblatt.sty 2026-10-12: Linie, Abstand, 13 pt, \pfjahre) ohne ihren Vorhalt von
% 14 Zeilen: der Setzer fordert vor dem Kopf 0,16 Seitenhöhe (Kopf + erste Aufgabe), mehr nicht
\setlength{\mbpfstufevorhalt}{0pt}
\makeatother
\newcommand{\karo}[1]{\par\smallskip\noindent\begin{tikzpicture}
  \clip (0,0) rectangle ({\linewidth-1mm},{#1*10mm});
  \draw[black!16,line width=0.3pt,step=5mm] (0,0) grid ({\linewidth},{#1*10mm});
  \end{tikzpicture}\par}
\newcommand{\kk}[1]{$\square$\,#1\hspace{0.9em}}
\newcommand{\linie}[1]{\,{\color{black!45}\rule[-1pt]{#1}{0.4pt}}\,}
\newcommand{\pkt}{\textperiodcentered{}}
% v0.14 (14.10.): Stufenkopf mit fünf Punkten (Option = Jahre für die Bestellkarte) oder „zuletzt 19“
\NewDocumentCommand{\gpunkte}{o m}{#2}
\newcommand{\gzuletzt}[1]{zuletzt #1}
% v0.7 (12.10.): Schülerfrage unter dem Stufenkopf; Ankreuzoption als Bild (Buchstabe und Kästchen darunter)
\newcommand{\gfrage}[1]{\par\nobreak\noindent\hspace*{\mbpfnr}{\small\itshape #1}\par\nobreak\addvspace{2pt}}
\newlength{\gbildbreite}\setlength{\gbildbreite}{0.32\linewidth}% Setzer setzt sie je Reihe (vier Bilder: 0,24)
\newcommand{\gbildoption}[3]{\begin{minipage}[t][#3][b]{\gbildbreite}\centering #2\par\smallskip\mbox{#1\,$\square$}\end{minipage}}
% Lineare (12.10.): Bild mit Buchstabe ohne Kästchen (Diagramm zuordnen); Aussage mit „wahr / falsch“ zum Ankreuzen
\newcommand{\gbildname}[3]{\begin{minipage}[t][#3][b]{\gbildbreite}\centering #2\par\smallskip\mbox{\textbf{#1}}\end{minipage}}
\newcommand{\gwahrfalsch}[2]{\par\noindent\hangindent2.2em\hangafter1\makebox[2.2em][l]{(#1)}#2\hfill\mbox{wahr $\square$\quad falsch $\square$}\par}
\begin{document}

\pfheftstil
\pflaufkopf[P10 \pkt{} Prüfung]{Netz und Schrägbild}{PG80-5}\pfseitevon
\pfheftkopf{Netz und Schrägbild}{}
\pfstufenzeile{Netz erkennen \pkt{} Netz skizzieren \pkt{} Schrägbild skizzieren}

% ===================== Stufe: Netz erkennen =====================
\Needspace*{0.16\textheight}\pfstufekopf{Netz erkennen}{\gpunkte[22]{\textcolor{black!25}{$\circ$}\,\textcolor{black!25}{$\circ$}\,\textcolor{black!25}{$\circ$}\,\textcolor{black!25}{$\circ$}\,\textcolor{black!80}{$\bullet$}}}
% Schülerfrage der Stufe (Gliederung; Einstieg nur von dort, nie aus der Bank – 12.10.)
\begin{pfaufg}{}{1.}{}
Stell es dir im Kopf vor und antworte, ohne zu rechnen.
\par\smallskip
\par\noindent\makebox[\mbpfteil][l]{\textbf{a)}}Drei Quadrate hängen in einer Reihe aneinander; du faltest sie zu einem Teil eines Würfels. Wo liegt dann das dritte Quadrat, vom ersten aus gesehen?\hfill\linie{30mm}
\par\noindent\makebox[\mbpfteil][l]{\textbf{b)}}Du rollst ein Rechteck zum Mantel eines Zylinders mit dem Durchmesser d. Wie lang ist die Seite, die einmal um den Kreis läuft?\hfill\linie{30mm}
\par\fusshilfe{\mbox{1:~a) gegenüber; b) $\pi$ · d, gut dreimal so lang wie d}}
\end{pfaufg}
% 2016-OS-B1j-m1 (Original 2016-OS-B1j, Schwere I, rang 6; herausgelöste Fassung: keine)
\begin{pfaufg}{’16~\pkt{}~1j}{2.}{Basis}
Aus dem abgebildeten Netz wird ein Würfel gefaltet. Nach dem Falten liegt die graue Fläche oben. Markiere die Fläche, die dann unten liegt.
\par\smallskip\noindent\begin{tikzpicture}[scale=0.8,line width=0.8pt]\fill[black!25] (-2,1) rectangle (-1,2);\draw (0,0) rectangle (1,3);\draw (0,1) -- (1,1) (0,2) -- (1,2);\draw (-2,1) rectangle (0,2);\draw (-1,1) -- (-1,2);\draw (1,0) rectangle (2,1);\end{tikzpicture}\par
\fusshilfe{\mbox{2:~das mittlere Quadrat der senkrechten Dreierreihe}}
\end{pfaufg}
% 2022-OS-K2a-m1 (Original 2022-OS-K2a, Schwere II, rang 3; herausgelöste Fassung: keine)
\begin{pfaufg}{nach~P10\\’22~\pkt{}~2a}{3.}{}
Ein Becher hat die Form eines Zylinders ohne Deckel. Er ist h = 7,5 cm hoch und hat den Radius r = 3,4 cm. Kreuze an, welches der drei Netze zum Becher passt. Ermittle die Länge und die Breite der Mantelfläche.
\par\smallskip\noindent\setlength{\gbildbreite}{\dimexpr0.96\linewidth/3\relax}\sbox{\gbox}{\vbox{\begin{tikzpicture}[line width=0.8pt]\draw (0,0) rectangle (2.72,1.5);\draw (3.4,0.75) circle (0.68);\end{tikzpicture}\par\begin{tikzpicture}[line width=0.8pt]\draw (0,0) rectangle (4.9,1.5);\draw (2.45,2.18) circle (0.68);\end{tikzpicture}\par\begin{tikzpicture}[line width=0.8pt]\draw (0,0) rectangle (4.27,1.5);\draw (3.59,2.18) circle (0.68);\end{tikzpicture}}}\gbildoption{A}{\begin{tikzpicture}[line width=0.8pt]\draw (0,0) rectangle (2.72,1.5);\draw (3.4,0.75) circle (0.68);\end{tikzpicture}}{\dimexpr\ht\gbox/3+\baselineskip\relax}\hfill\gbildoption{B}{\begin{tikzpicture}[line width=0.8pt]\draw (0,0) rectangle (4.9,1.5);\draw (2.45,2.18) circle (0.68);\end{tikzpicture}}{\dimexpr\ht\gbox/3+\baselineskip\relax}\hfill\gbildoption{C}{\begin{tikzpicture}[line width=0.8pt]\draw (0,0) rectangle (4.27,1.5);\draw (3.59,2.18) circle (0.68);\end{tikzpicture}}{\dimexpr\ht\gbox/3+\baselineskip\relax}\par
\karo{3}
\fusshilfe{\mbox{3:~C; 21,4 cm; 7,5 cm}}
\end{pfaufg}

% ===================== Stufe: Netz mit Maßen skizzieren =====================
% Formel-Frage (einmal je Blatt, erste Stufe mit Formel; pruefung.md § 5)
\Needspace*{0.16\textheight}\pfstufekopf{Netz skizzieren}{\gpunkte[23 22]{\textcolor{black!25}{$\circ$}\,\textcolor{black!25}{$\circ$}\,\textcolor{black!25}{$\circ$}\,\textcolor{black!80}{$\bullet$}\,\textcolor{black!80}{$\bullet$}}}\gfrage{Bevor du rechnest: Auf welcher Seite der Formelsammlung steht die Formel? Seite~\linie{12mm}}
% Schülerfrage der Stufe (Gliederung; Einstieg nur von dort, nie aus der Bank – 12.10.)
\begin{pfaufg}{}{4.}{}
Ein Prisma hat als Grundfläche ein gleichseitiges Dreieck mit 2 cm Seitenlänge und ist 5 cm hoch. Im Netz hängt ein Rechteck des Mantels an einer Seite des Dreiecks.
\par\smallskip
\par\noindent Wie lang ist die Seite des Rechtecks, die am Dreieck anliegt, und wie lang die andere?\hfill\linie{30mm}
\par\fusshilfe{\mbox{4:~2 cm und 5 cm}}
\end{pfaufg}
% 2023-OS-K5a-m1 (Original 2023-OS-K5a, Schwere II, rang 2; herausgelöste Fassung: keine)
\begin{pfaufg}{nach~P10\\’23~\pkt{}~5a}{5.}{}
Eine Konservendose hat die Form eines Zylinders mit dem Radius 3,5 cm und der Höhe 7 cm. Skizziere die Mantelfläche der Dose. Bestimme ihre Länge a und ihre Breite b und schreibe beide an deine Skizze.
\gzwei{\zylinder{1.1}{2.4}{r}{h}}{\karo{4}}
\fusshilfe{\mbox{5:~22,0 cm; 7 cm}}
\end{pfaufg}
% 2019-OS-K4a-m1 (Original 2019-OS-K4a, Schwere II, rang 9; herausgelöste Fassung: keine)
\begin{pfaufg}{nach~P10\\’19~\pkt{}~4a}{6.}{}
Auf einem Mast sitzt ein dreiseitiges Prisma. Seine drei Seitenflächen sind Rechtecke mit 1,80 m Breite und 1,50 m Höhe; Grund- und Deckfläche sind gleichseitige Dreiecke. Das Netz des Prismas ist schon begonnen. Vervollständige das Netz. Gib die Länge a an.
\par\smallskip\noindent\begin{tikzpicture}[scale=0.5,line width=0.8pt]\draw[black!14,line width=0.3pt] (0,0) grid (19,16);\draw (0.5,5.5) rectangle (6.5,10.5);\draw (6.5,5.5) rectangle (12.5,10.5);\draw (6.5,5.5) -- (12.5,5.5) -- (9.5,0.304) -- cycle;\node[font=\small,right,inner sep=1.5pt] at (11.1,2.7) {$a$};\node[font=\scriptsize,mbgrau,below right,inner sep=1pt] at (0,0) {1 Kästchen $\mathrel{\widehat{=}}$ 0,30 m};\end{tikzpicture}\par
\par\noindent\hfill a $=$\linie{30mm}
\fusshilfe{\mbox{6:~drittes Rechteck 6 × 5 Kästchen an das rechte, zweites Dreieck oben an ein Rechteck; 1,80 m}}
\end{pfaufg}

% ===================== Stufe: Körper im Schrägbild skizzieren =====================
\Needspace*{0.16\textheight}\pfstufekopf{Schrägbild skizzieren}{\gpunkte[24]{\textcolor{black!25}{$\circ$}\,\textcolor{black!25}{$\circ$}\,\textcolor{black!80}{$\bullet$}\,\textcolor{black!25}{$\circ$}\,\textcolor{black!25}{$\circ$}}}
% Schülerfrage der Stufe (Gliederung; Einstieg nur von dort, nie aus der Bank – 12.10.)
\begin{pfaufg}{}{7.}{}
Ein Kegel mit dem Radius 3 cm und der Höhe 8 cm steht aufrecht in einer Schachtel.
\par\smallskip
\par\noindent\makebox[\mbpfteil][l]{\textbf{a)}}Wie breit muss die Schachtel mindestens sein?\hfill\linie{30mm}
\par\noindent\makebox[\mbpfteil][l]{\textbf{b)}}Wie zeichnest du den Grundkreis im Schrägbild?\hfill\linie{30mm}
\par\fusshilfe{\mbox{7:~a) 6 cm, so breit wie der Durchmesser; b) als flache Ellipse}}
\end{pfaufg}
% 2024-OS-K4b-m1 (Original 2024-OS-K4b, Schwere II, rang 1; herausgelöste Fassung: keine)
\begin{pfaufg}{nach~P10\\’24~\pkt{}~4b}{8.}{}
Eine Firma stellt Kegel mit dem Radius r = 25 cm und der Höhe h = 60 cm her. Der Kegel soll in einem quaderförmigen Karton verschickt werden, der möglichst klein ist, in den der Kegel aber ganz hineinpasst. Skizziere den Kegel in den Karton. Kreuze die kleinste passende Größe an (Länge × Breite × Höhe):
\gzwei{\quader{2.2}{2.2}{4.4}{}{}{}}{\vspace{2mm}\pfkreuzab{A}{26 cm × 26 cm × 61 cm}\pfkreuzab{B}{51 cm × 51 cm × 81 cm}\pfkreuzab{C}{51 cm × 51 cm × 61 cm}\pfkreuzab{D}{26 cm × 51 cm × 61 cm}}
\par\smallskip\noindent Begründe deine Wahl.
\karo{2}
\fusshilfe{\mbox{8:~Kegel steht auf dem Boden, Grundkreis als Ellipse, Spitze mittig oben; C; je 1 cm mehr als Durchmesser 50 cm und Höhe 60 cm – 26 cm sind schmaler als der Durchmesser, 81 cm sind unnötig hoch}}
\end{pfaufg}

\pfblattende{11}{26 & 25 & 24 & 23 & 22 & 21 & 20 & 19 & 18 & 17 & 16}{ &  & 4b & 5a & 2a &  & 5a & 4a & 1i & 1c & 1j}%
{\pfausschnitt{Volumen und Oberfläche}{Netz und Schrägbild}{}}
\end{document}